\section*{Chapter \ref{ch:g9:acids-bases-ph} --- Acids, Bases and pH}

\begin{solution}{exo:g9:acids-bases-ph:1}
\omterm{def:g9:acids-bases-ph:ph}{pH} 3: acidic. \omterm{def:g9:acids-bases-ph:ph}{pH} 7: neutral. \omterm{def:g9:acids-bases-ph:ph}{pH} 11: basic.
\end{solution}

\begin{solution}{exo:g9:acids-bases-ph:2}
Acidic: \ce{H+} \omterm{def:g9:ions:ion}{ions}. Basic: \ce{OH-} \omterm{def:g9:ions:ion}{ions}.
\end{solution}

\begin{solution}{exo:g9:acids-bases-ph:3}
Lime juice (2), coffee (5), pure water (7), sea water (8.1), household
ammonia (12).
\end{solution}

\begin{solution}{exo:g9:acids-bases-ph:4}
A strip on a clean dry dish is touched with a drop of the solution
carried by a clean glass rod; the colour is compared at once with the
chart.
\end{solution}

\begin{solution}{exo:g9:acids-bases-ph:5}
About 5 after a tenfold dilution; about 6 after a hundredfold one.
\end{solution}

\begin{solution}{exo:g9:acids-bases-ph:6}
No. Diluting brings the \omterm{def:g9:acids-bases-ph:ph}{pH} closer to 7, but never past it: water itself
is neutral, and adding it cannot make \ce{OH-} \omterm{def:g9:ions:ion}{ions} outnumber \ce{H+}
\omterm{def:g9:ions:ion}{ions}.
\end{solution}

\begin{solution}{exo:g9:acids-bases-ph:7}
Mixing gives out heat. Water poured onto a concentrated acid could boil
at once and throw acid out; acid poured slowly into water is spread at
once in a large volume.
\end{solution}

\begin{solution}{exo:g9:acids-bases-ph:8}
Corrosive: burns skin and eyes, attacks metals. Wear goggles and gloves;
rinse any splash with plenty of water.
\end{solution}

\begin{solution}{exo:g9:acids-bases-ph:9}
Three dilutions (1 to 2, 2 to 3, 3 to 4): $10 \times 10 \times 10 =
1000$ times.
\end{solution}

\begin{solution}{exo:g9:acids-bases-ph:10}
Milk of magnesia is basic (\omterm{def:g9:acids-bases-ph:ph}{pH} about 10.5): it reacts with part of the
excess acid of the stomach.
\end{solution}

\begin{solution}{exo:g9:acids-bases-ph:11}
No: with a \omterm{def:g9:acids-bases-ph:ph}{pH} of 8.1 it is still slightly basic. It is moving towards
lower \omterm{def:g9:acids-bases-ph:ph}{pH}, less basic, because carbon dioxide from the \omterm{def:g4:air-a-mixture-of-gases:air}{air} \omterm{def:g2:mixing-and-dissolving:dissolve}{dissolves} in
it.
\end{solution}

\begin{solution}{exo:g9:acids-bases-ph:12}
\omterm{def:g9:acids-bases-ph:ph-paper}{pH paper} is read by eye against a colour chart, to about one unit; the
\omterm{def:g9:acids-bases-ph:ph-paper}{pH meter} gives a number to one or two decimals. The meter is more
precise; the two agree within the precision of the paper.
\end{solution}

\begin{solution}{pb:g9:acids-bases-ph:1}
\textbf{1.} Acidic: it contains mostly \ce{H+(aq)} and \ce{Cl-(aq)}
\omterm{def:g9:ions:ion}{ions}.

\textbf{2.} GHS05, corrosive: an acid attacks skin, eyes and metals.

\textbf{3.} It gives a precise number instead of a colour to compare by
eye.

\textbf{4.} Ten times more: one \omterm{def:g9:acids-bases-ph:ph}{pH} unit corresponds to a tenfold dilution.

\textbf{5.} \qty{10}{L}, of \omterm{def:g9:acids-bases-ph:ph}{pH} 3.

\textbf{6.} Three: \omterm{def:g9:acids-bases-ph:ph}{pH} 2 to 3, 3 to 4, 4 to 5.

\textbf{7.} $1 \times 10 \times 10 \times 10 = \qty{1000}{L}$.

\textbf{8.} No: diluting only brings the \omterm{def:g9:acids-bases-ph:ph}{pH} closer to 7.

\textbf{9.} \ce{H+ + OH- -> H2O}.

\textbf{10.} The acid is destroyed, turned into water, instead of being
spread through a huge volume; only a little base is needed and no
thousand litres of water.

\textbf{11.} Baking soda dissolved in water (or soapy water).

\textbf{12.} $1000 - 1 = \qty{999}{L}$ of water.
\end{solution}
